\documentclass[11pt]{article}
\usepackage[T1]{fontenc}
\usepackage{lmodern,microtype,amsmath,amssymb,amsthm,booktabs}
\usepackage[margin=1in]{geometry}
\usepackage{xurl}
\usepackage[hidelinks]{hyperref}
\newtheorem{theorem}{Theorem}
\newtheorem{lemma}{Lemma}
\newcommand{\R}{\mathbb R}
\newcommand{\dd}{\Delta}
\title{An explicit finite-kinase instability bound\protect\\
for a mixed-mechanism phosphorylation system}
\author{Zijian Zeng}
\date{}
\hypersetup{pdftitle={An explicit finite-kinase instability bound for a mixed-mechanism phosphorylation system},pdfauthor={Zijian Zeng}}
\begin{document}
\maketitle
\begin{abstract}
For the mass-action mixed-mechanism phosphorylation system of Conradi,
Mincheva and Shiu, at their Table~1 rates and fixed phosphatase and substrate
totals $8$ and $40$, we prove that the unique positive equilibrium has exactly
two unstable and four stable compatibility-class eigenvalues whenever the
kinase total is at least $4660166/209569<23$. A rational equilibrium
parameterization converts the unbounded range of kinase totals into a bounded
interval. Exact Bernstein certificates for six Hurwitz determinants then
give a uniform Routh root count. The bound is sufficient; sharpness is not
asserted. This proves the first assertion of their Conjecture~6.2 with explicit
witness totals. Its periodic-orbit assertion, and hence the full original
problem, remain unresolved here.
\end{abstract}

\section{Model, statement and relation to prior work}

Consider processive phosphorylation and distributive dephosphorylation:
\begin{align*}
 S_0+K&\mathrel{\mathop{\rightleftarrows}^{k_1}_{k_2}}S_0K
       \xrightarrow{k_3}S_1K\xrightarrow{k_4}S_2+K,\\
 S_2+P&\mathrel{\mathop{\rightleftarrows}^{k_5}_{k_6}}S_2P
       \xrightarrow{k_7}S_1+P
       \mathrel{\mathop{\rightleftarrows}^{k_8}_{k_9}}S_1P
       \xrightarrow{k_{10}}S_0+P.
\end{align*}
Use the species order
\[
 (x_1,\ldots,x_9)=(S_0,K,S_0K,S_1K,S_2,P,S_2P,S_1,S_1P)
\]
and the fixed rate vector
\begin{equation}\label{eq:rates}
 (k_1,\ldots,k_{10})=(1,1,1,1,100,1,9/10,3,1,100).
\end{equation}
Thus, in the notation of Conradi--Mincheva--Shiu~\cite{CMS}, the equations are
\begin{equation}\label{eq:odes}
\begin{aligned}
 \dot x_1&=-x_1x_2+x_3+100x_9,&
 \dot x_2&=-x_1x_2+x_3+x_4,\\
 \dot x_3&=x_1x_2-2x_3,& \dot x_4&=x_3-x_4,\\
 \dot x_5&=x_4-100x_5x_6+x_7,&
 \dot x_7&=100x_5x_6-\tfrac{19}{10}x_7,\\
 \dot x_6&=-100x_5x_6-3x_8x_6+\tfrac{19}{10}x_7+101x_9,\\
 \dot x_8&=\tfrac9{10}x_7-3x_6x_8+x_9,&
 \dot x_9&=3x_6x_8-101x_9.
\end{aligned}
\end{equation}
The conserved totals are
\begin{equation}\label{eq:totals}
 \kappa=x_2+x_3+x_4,\qquad P_{\rm tot}=x_6+x_7+x_9,\qquad
 S_{\rm tot}=x_1+x_3+x_4+x_5+x_7+x_8+x_9.
\end{equation}
Here $\kappa$ denotes the total kinase, to distinguish it from free kinase.
The three independent conservation laws define a six-dimensional
compatibility class in the nonnegative orthant.

\begin{theorem}\label{thm:main}
For \eqref{eq:odes} with $P_{\rm tot}=8$ and $S_{\rm tot}=40$, and every finite
\[
 \kappa\ge B:=\frac{4660166}{209569},
\]
the unique positive equilibrium has exactly two eigenvalues in
$\{\lambda:\operatorname{Re}\lambda>0\}$ and four in
$\{\lambda:\operatorname{Re}\lambda<0\}$ on its compatibility class,
counted with algebraic multiplicity. In particular, this holds for every
$\kappa\ge23$.
\end{theorem}

The source paper~\cite{CMS}, especially Proposition~3.1, Proposition~4.1 and
Section~6, provides equilibrium uniqueness, a stability-surface analysis and
large-kinase simulations. Its Conjecture~6.2(1) asks for fixed phosphatase and
substrate totals admitting unstable equilibria at arbitrarily large kinase
totals. Theorem~\ref{thm:main} supplies these witnesses and certifies an entire
unbounded ray, excluding restabilization on that ray. Equilibrium
parameterization and the Hurwitz method themselves are established tools.

There is direct subsequent topical overlap: Grasse's 2023 poster-programme
entry concerns kinase-total robustness~\cite{Grasse2023}, and the 2025 abstract
of Grasse and Conradi considers the reversed processive/distributive
arrangement~\cite{GC2025}. The inspected programme and abstract give no
equivalent finite bound or complete proof. The poster itself and possible
follow-up texts were unavailable to this comparison. No priority claim is
made. The contribution here is the explicit bound and exact certificate for
the stated fixed totals, rather than the prior observation of oscillations.

\section{Equilibrium parameter and coverage of the kinase ray}

Put $z=x_9$ and define
\begin{equation}\label{eq:param}
\begin{aligned}
 p(z)&=8-\frac{1009}{9}z,&
 G(z)&=\frac{2809}{9}z+\frac{322z}{9p(z)},\\
 r(z)&=40-G(z),&
 \kappa(z)&=200z+\frac{200z}{r(z)}.
\end{aligned}
\end{equation}
At a positive equilibrium, direct balance in \eqref{eq:odes} yields
\begin{equation}\label{eq:eq}
 x^*(z)=\left(r,\frac{200z}{r},100z,100z,
       \frac{19z}{9p},p,\frac{1000z}{9},\frac{101z}{3p},z\right).
\end{equation}
For example, $\dot x_4=\dot x_1=\dot x_3=0$ gives
$x_4=x_3=100z$ and $x_1x_2=200z$.
The $x_5,x_7$ balances give $x_7=1000z/9$ and
$x_5=19z/(9p)$; the $x_9$ balance gives $x_8=101z/(3p)$.
The phosphatase and substrate totals then give $p$ and $r$.
Conversely, substitution of \eqref{eq:eq} into all nine equations and
\eqref{eq:totals} verifies all these identities.

On $0<z<72/1009$ we have $p>0$ and
\[
 G'(z)=\frac{2809}{9}+\frac{2576}{9p(z)^2}>0.
\]
Since $G(0)=0$ and $G(z)\to\infty$ at the right endpoint, there is exactly
one $\alpha$ in this interval satisfying $G(\alpha)=40$.
The positive equilibria correspond precisely to $0<z<\alpha$.
On this interval,
\[
 \kappa'(z)=200+\frac{200}{r(z)}
                 +\frac{200zG'(z)}{r(z)^2}>0.
\]
The endpoint limits of $\kappa(z)$ are $0$ and $\infty$.
Thus \eqref{eq:eq} gives exactly one positive equilibrium for every positive
finite total kinase. No interval of totals is omitted.

For an explicit enclosure define
\begin{equation}\label{eq:interval}
 Q(z)=2834281z^2-568386z+25920,\qquad
 L=\frac7{100},\qquad U=\frac{70122383355}{10^{12}}.
\end{equation}
Then $r=Q/[9(72-1009z)]$, and exact evaluation gives
\[
 Q(L)=\frac{209569}{10000}>0,\qquad
 Q(U)=-\frac{6175117557936279}{40000000000000000000000}<0.
\]
Moreover $0<L<U<72/1009$ and $Q'(U)<0$.
Because $Q''>0$, the polynomial is strictly decreasing on $[0,U]$.
Consequently $L<\alpha<U$, and $Q>0$ on $[L,\alpha)$.
Finally $\kappa(L)=B<23$. Strict monotonicity proves that
$[L,\alpha)$ maps bijectively onto $[B,\infty)$.

\section{Exact signs and eigenvalue count}

In the compatibility class retain coordinates
$(x_3,x_4,x_5,x_7,x_8,x_9)$ and eliminate
\[
 x_1=40-x_3-x_4-x_5-x_7-x_8-x_9,\quad
 x_2=\kappa-x_3-x_4,\quad x_6=8-x_7-x_9.
\]
Differentiate at fixed $\kappa$, then substitute \eqref{eq:eq} and
$\kappa=\kappa(z)$. With $u=200z/r$, $w=19z/(9p)$ and $v=101z/(3p)$,
the resulting Jacobian is
\begin{equation}\label{eq:jac}
 J(z)=\begin{pmatrix}
 -u-r-2&-u-r&-u&-u&-u&-u\\
 1&-1&0&0&0&0\\
 0&1&-100p&100w+1&0&100w\\
 0&0&100p&-100w-19/10&0&-100w\\
 0&0&0&9/10+3v&-3p&1+3v\\
 0&0&0&-3v&3p&-3v-101
 \end{pmatrix}.
\end{equation}
This is a restriction to the tangent space of the conservation laws, so its
six eigenvalues exclude the three conservation zeros of the full system.
In particular, $\kappa(z)$ is not differentiated as a function of state.

Write
\[
 \det(\lambda I-J(z))=\sum_{i=0}^{6}c_i(z)\lambda^{6-i},\qquad c_0=1,
 \qquad
 \dd_j(z)=\det\big[c_{2b-a+1}(z)\big]_{a,b=0}^{j-1},
\]
where $c_i=0$ outside $0\le i\le6$. This is the usual Hurwitz convention;
its transpose gives the same leading determinants.
Let $A=1009z-72$ and use the denominators $D_j$ in Table~\ref{tab:signs}.
Define the polynomials $N_j=D_j\dd_j$ and
\[
 h_j=-N_j\quad(1\le j\le5),\qquad h_6=N_6.
\]
The displayed Jacobian and these definitions specify every polynomial without
the need for expanded, large integer expressions.

\begin{lemma}\label{lem:sign}
For all $z\in[L,\alpha)$,
\[
 (\operatorname{sgn}\dd_1,\ldots,\operatorname{sgn}\dd_6)
       =(+,+,+,+,-,-).
\]
\end{lemma}

\begin{proof}
For a polynomial $h$ of degree $n$, put
\[
 a_\ell=\frac{h^{(\ell)}(L)}{\ell!}(U-L)^\ell,\qquad
 b_k=\sum_{\ell=0}^{k}a_\ell
                  \frac{\binom{k}{\ell}}{\binom{n}{\ell}}.
\]
The exact change of basis is
\begin{equation}\label{eq:bernstein}
 h(L+(U-L)t)=\sum_{k=0}^{n}b_k\binom nk t^k(1-t)^{n-k}.
\end{equation}
The basis functions on the right are nonnegative and sum to one on $[0,1]$.
Strict positivity of every $b_k$ therefore proves strict positivity of $h$
on the closed interval $[L,U]$.

Exact rational arithmetic applied to \eqref{eq:jac} gives
Table~\ref{tab:signs}. Its last two columns are conservative rational lower
bounds, not rounded estimates: for each row, $b_0=h_j(L)>10^{m_j}$ and
$b_k/b_0>\eta_j$ for every $0\le k\le n_j$. Thus all
$5+11+18+25+31+35=125$ Bernstein coefficients are strictly positive.
The supplementary verifier reconstructs these calculations from
\eqref{eq:jac}; expanded polynomials and all rational coefficients accompany
the note.

\begin{table}[htbp]
\centering
\small
\begin{tabular}{@{}clrrrr@{}}
\toprule
$j$& $D_j$ & $\operatorname{sgn}D_j$ & $n_j$ & $m_j$ & $\eta_j$\\
\midrule
1&$90AQ$&$-$&4&5&$24/1000$\\
2&$72900A^3Q^2$&$-$&10&14&$28/10000$\\
3&$177147000A^5Q^3$&$-$&17&24&$14/10000$\\
4&$71744535000A^7Q^4$&$-$&24&34&$8/10000$\\
5&$871696100250A^8Q^5$&$+$&30&41&$4/10000$\\
6&$94143178827A^7Q^6$&$-$&34&47&$4/10000$\\
\bottomrule
\end{tabular}
\caption{Exact sign certificate. Denominator signs are used only on
$[L,\alpha)$; numerator positivity is certified on $[L,U]$. The exponent
$m_j$ means $b_0>10^{m_j}$, and $\eta_j$ is a strict lower bound for every
$b_k/b_0$.}\label{tab:signs}
\end{table}

On the physical interval, $A<0$ and $Q>0$, so the denominator signs in the
table follow directly. The signs of $N_1,\ldots,N_6$ are
$(-,-,-,-,-,+)$ by \eqref{eq:bernstein}. Division gives the claimed signs.
Only the numerator polynomials are certified beyond $\alpha$: no determinant
is evaluated across its pole there.
\end{proof}

\begin{proof}[Proof of Theorem~\ref{thm:main}]
The Routh array is regular because every leading Hurwitz determinant is
nonzero. Its first column is
\[
 1,\quad \dd_1,\quad\frac{\dd_2}{\dd_1},\quad
 \frac{\dd_3}{\dd_2},\quad\frac{\dd_4}{\dd_3},\quad
 \frac{\dd_5}{\dd_4},\quad\frac{\dd_6}{\dd_5}.
\]
By Lemma~\ref{lem:sign}, its signs are $(+,+,+,+,+,-,+)$, with exactly
two sign changes. The classical Routh root-count theorem gives two roots in
the open right half-plane, counted with algebraic multiplicity.
There are no imaginary-axis roots: $c_6=0$ would force $\dd_6=0$;
a nonzero imaginary conjugate pair would give a common root of the even
and odd parts of the characteristic polynomial, forcing $\dd_5=0$.
Both alternatives are excluded. Thus the other four roots lie in the open
left half-plane. The bijection $[L,\alpha)\to[B,\infty)$ proves the result
for every total stated in the theorem, including $\kappa=B$.
\end{proof}

\section{Reproducibility and limitations}

The supplementary file \texttt{verify\_certificate.py} constructs
\eqref{eq:jac}, computes the characteristic polynomial and the six Hurwitz
determinants over rational functions, and checks every bound in
Table~\ref{tab:signs}. Specifically, it checks the displayed denominator
factorizations and degrees, computes the derivatives defining $a_\ell$,
forms every $b_k$, and compares them with the exact rational lower bounds.
It also checks its reconstructed characteristic polynomial and all 125
coefficients against the supplied full certificates.
All sign decisions are exact; no integration, fitted asymptotic law,
floating-point eigenvalue calculation or exchange of limits is used.
The calculations were performed with Python and SymPy~1.13.3. A separate
verification pass reconstructed the published nine equations in a different
affine coordinate chart and checked the complete polynomial identities using
rational arithmetic, including a standard-library \texttt{Fraction}
expansion of the Bernstein certificates.

The theorem is restricted to \eqref{eq:rates} and the fixed totals $8,40$.
The threshold is sufficient and is not asserted to be the stability boundary.
An unstable equilibrium in six dimensions does not by itself establish a
periodic orbit. Neither periodic-orbit existence for all sufficiently large
kinase totals nor invariance of its $S_2$ projection is proved here. These are
the additional requirements of Conjecture~6.2(2) in~\cite{CMS}; the full
original problem remains unresolved by this note.

\paragraph{Computational and AI assistance.}
OpenAI Codex assisted the symbolic derivation, exact-verification code,
literature screening, drafting and editing. A separate model context checked
the frozen mathematical argument and independently reconstructed its algebra.
This assistance and the executable checks are not formal proof-assistant
verification, human peer review or a guarantee of priority. The mathematical
argument and exact certificate specify the claims supported here.

\begin{thebibliography}{9}
\bibitem{CMS}
C.~Conradi, M.~Mincheva and A.~Shiu,
Emergence of oscillations in a mixed-mechanism phosphorylation system,
\emph{Bulletin of Mathematical Biology} \textbf{81} (2019), 1829--1852.
\href{https://doi.org/10.1007/s11538-019-00580-6}{doi:10.1007/s11538-019-00580-6}.
Model and conjecture numbering follow
\href{https://arxiv.org/abs/1809.02886v2}{arXiv:1809.02886v2}.

\bibitem{Grasse2023}
K.-J.~Grasse,
Analysis of the robustness of oscillations in a mixed mechanism double
phosphorylation system with respect to total kinase concentrations,
\emph{Systems Biology of Human Disease 2023}, poster programme,
poster~48, p.~52 (2023).
\href{https://www.sbhdberlin.org/pdf/SBHD2023_Poster_Session_Program.pdf}{Official poster programme}.
Programme entry only; the poster text was not inspected.

\bibitem{GC2025}
K.-J.~Grasse and C.~Conradi,
On the robustness of oscillations in a mixed mechanism phosphorylation system
against perturbations in the total amount of the processive enzyme,
\emph{BIOMATH 2025 Conference Abstracts}, Sofia, 15--20 June 2025.
\href{https://biomath.math.bas.bg/biomath/index.php/bmcs/article/download/1787/1594/7828}{Official abstract}.
Conference abstract.
\end{thebibliography}
\end{document}
